%% kapa: exam=1 round=25 year=2014 period=1 form=A subject=economics subject_ko=경제원론 seq=31 number=71
%% source: Q-Net 감정평가사 기출문제 (cst00309) artlSeq=5101207 file=extracted/1008/5101207/2014년 제25회 감정평가사 1차시험 기출문제/1교시(A형).hwp
%% conversion: kordoc 4.24.51 -> tools/gen1.py (auto); review=proofread-sonnet; flags=none
\begin{question}{71}
소규모 개방국가인 A국과 B국의 통화량 증가율은 매년 각각 5 \%와 3 \%이다. 두 국가의 실질 GDP 증가율은 매년 2 \%로 일정하고 여타 면에서도 서로 동일하다. 이 때 두 국가의 장기균형에 관한 설명으로 옳지 않은 것은? (단, 두 국가의 명목환율은 A국 통화 1단위와 교환되는 B국 통화의 양으로 정의한다.)
\begin{choices}
  \item 명목환율은 하락할 것이다.
  \item A국의 물가상승률이 B국보다 더 높을 것이다.
  \item B국의 명목이자율이 A국보다 더 낮을 것이다.
  \item A국의 명목GDP 성장률이 B국보다 더 높을 것이다.
  \item A국은 무역수지 흑자, B국은 무역수지 적자가 발생할 것이다.
\end{choices}
\end{question}
\answer{5}
